\section{Continual Learning, World Model y Expert Agent}

Cuando el conductor preguntó cuál es el mayor cuello de botella de los Agent, Su Yu (苏煜) presentó memory, self learning, continual learning, world model, specialization y expert agent como distintas facetas de una misma cuestión. Lo que hoy les falta a los Agent es aprender de forma continua a partir de la interacción y convertir la experiencia en capacidades estables.

\begin{figure}[H]
\centering
\includegraphics[width=0.94\textwidth]{continual-learning-map.png}
\caption{Relación entre Continual Learning, World Model y Expert Agent.}
\end{figure}

\begin{knowledgebox}{Lectura de la figura: por qué estos términos forman una cadena}
La experiencia de la izquierda proviene de la interacción con el entorno; el continual/self learning actualiza esa experiencia y la convierte en capacidades; el world model representa el aprendizaje de las regularidades del entorno; la specialization representa la formación de un experto de dominio; el resultado final es un Agent más confiable, más rápido y de menor costo. La conclusión que sustenta es que memory, world model y la especialización no son roadmaps independientes entre sí, sino una misma cadena de acumulación de capacidades.
\end{knowledgebox}

\subsection{Términos clave: vocabulario de los cuellos de botella de los Agent}

\noindent\begin{tabularx}{\textwidth}{p{0.20\textwidth}p{0.32\textwidth}X}
\toprule
Término & Problema que resuelve & Lugar en este episodio \\
\midrule
Memory & Conservar el historial y las preferencias relevantes en tareas de varios pasos & Sin memoria, el Agent actúa cada vez como novato y tiende a repetir errores. \\
Continual Learning & Actualizar las capacidades de forma continua a partir de nuevas interacciones & Hace que el Agent no solo ejecute guiones, sino que acumule experiencia. \\
World Model & Aprender el estado del entorno, la causalidad y qué acciones son posibles & Sustenta la planificación, la predicción de consecuencias y la operación segura. \\
Specialization & Convertirse en expert agent de un dominio & Pasar de asistente general a trabajador vertical confiable. \\
Reliability & Tasa de finalización estable de las tareas & Uno de los mayores problemas actuales en la conversión en producto de los Agent. \\
Cost Effectiveness & Si el costo por tarea es aceptable & Determina si el Agent puede pasar de la demo a producción. \\
\bottomrule
\end{tabularx}

\begin{importantbox}{Formulación unificada del mayor cuello de botella}
El mayor cuello de botella de los Agent no es un módulo de memory aislado, un prompt aislado ni una herramienta aislada, sino la incapacidad de transformar, como lo hace una persona, la experiencia de largo plazo en capacidades profesionales estables. Continual learning, world model y specialization sirven, en última instancia, a la confiabilidad, la velocidad y el costo.
\end{importantbox}

\subsection{La señal de los forward-deployed engineers}

En la entrevista se menciona que algunas empresas recurren a forward-deployed engineers: envían ingenieros a las instalaciones del cliente para ayudarle a construir sus Agents. Esto indica que los Agent todavía no se han difundido con una barrera de entrada baja: hace falta entender los procesos del cliente, el entorno de herramientas, los modos de falla y las restricciones de seguridad. El producto sigue en una etapa con un alto componente de servicio.

\begin{warningbox}{La dificultad del despliegue no se debe solo a un modelo débil}
Que a los Agent les cueste entrar en las empresas no se debe solo a la insuficiencia de capacidades del modelo, sino también a la comprensión de los procesos, la gestión de permisos, la integración de datos, la responsabilidad por los errores, los límites de seguridad y el cálculo de costos. Atribuir todas las fallas a que «el modelo no es lo bastante inteligente» lleva a juzgar mal la dificultad de la conversión en producto.
\end{warningbox}

\subsection{Resumen de la sección}

La línea principal de la siguiente etapa de los Agent es el aprendizaje continuo y el modelado del mundo. El progreso real se manifestará en confiabilidad, rapidez, bajo costo y especialización, y no solo en completar secuencias de acciones más vistosas en una demo.
